Let $BD=CD=t$ and $BC=a$. The angle-bisector theorem gives $\dfrac{8}{t}=\dfrac{12}{a}$, so $a=\dfrac{3t}{2}$. The angle-bisector length identity is $BD^2=AB\cdot BC-AD\cdot DC$, obtainable directly from Stewart' s theorem and the bisector ratio. Hence $t^2=12(\dfrac{3t}{2})-8t=10t$, giving $t=10$ and $\boxed{BC=15}$.
\begin{asy}
import graph;import geometry;usepackage("amsmath");size(7cm);
pair A=(0,0),C=(18,0),D=(8,0),B=(27/4,15sqrt(7)/4);draw(A--B--C--cycle);draw(B--D,blue);label("$A$",A,(-1,-1));label("$B$",B,(0,1));label("$C$",C,(1,-1));label("$D$",D,(0,-1));label("$12$",(A+B)/2,(-1,1));label("$8$",(A+D)/2,(0,-1));label("$t$",(B+D)/2,(1,0));label("$t$",(C+D)/2,(0,-1));
\end{asy}

